En este apéndice se recopilan de forma íntegra y literal las respuestas de los usuarios que participaron en las pruebas de validación de \textit{Reset} a las dos preguntas de respuesta abierta planteadas en el cuestionario de experiencia de usuario. 

La recopilación de estas opiniones permite analizar en detalle las fortalezas del videojuego y los aspectos susceptibles de mejora o depuración técnica identificados por los evaluadores.

% Incrementar la altura de las celdas para que tengan más espacio vertical (padding)
\renewcommand{\arraystretch}{1.4}

\section{Aspectos más valorados (Fortalezas)}
La Tabla~\ref{tab:comentarios_mas_gustado} recoge las opiniones de los usuarios al ser preguntados por el aspecto que más les gustó del videojuego.

\begin{longtable}{|l|p{12.5cm}|}
\hline
\textbf{Usuario} & \textbf{¿Qué es lo que MÁS te ha gustado de Reset?} \\ \hline
\endfirsthead
\hline
\textbf{Usuario} & \textbf{¿Qué es lo que MÁS te ha gustado de Reset? (Continuación)} \\ \hline
\endhead
\hline
\endfoot
\hline
\endlastfoot
Usuario 1 & Me ha hecho aprender rápidamente a jugar en ordenador, cosa que ni había hecho nunca. \\ \hline
Usuario 2 & Lo que más me han gustado han sido las interfaces de los mundos y la ambientación general. \\ \hline
Usuario 3 & Que hace que tengas que probar varios tipos de juego, además que habitualmente no son de los que más juego. \\ \hline
Usuario 4 & La historia, la ambientación y las mecánicas. \\ \hline
Usuario 5 & Las referencias a otros grandes videojuegos de la industria. \\ \hline
Usuario 6 & Lo bien que se siente la cohesión de la estética de los niveles. \\ \hline
Usuario 7 & Resolver el puzle en \textit{The Last Reset 2}. \\ \hline
Usuario 8 & Los menús y la UI están genial. \\ \hline
Usuario 9 & A pesar de no estar nada acostumbrado a jugar en teclado, me ha parecido que el personaje en \textit{Super Reset Bros} se controlaba muy bien. Los enemigos del Nivel 1 estaban bien situados y, si se dominan los controles, se pueden hacer saltos increíbles rebotando de un enemigo a otro, lo cual incrementa la satisfacción del jugador. En cuanto al Nivel 2 del mismo mundo, es muy satisfactorio impulsarte por las paredes rápidamente y utilizar el \textit{dash} para realizar saltos muy largos y salvar varias zonas de pinchos. Siento que \textit{Super Reset Bros} está muy pulido. \\ \hline
Usuario 10 & La jugabilidad y fluidez del juego con sus respectivas mecánicas y sobre todo destacar la mezcla de estéticas según cada juego pero manteniendo una coherencia global que te hace sentir que estás en un mismo juego. \\ \hline
Usuario 11 & La mezcla de diferentes géneros de videojuegos en uno solo con sus diferentes estéticas y controles. \\ \hline
Usuario 12 & Lo fluido que funcionan todos los controles y lo bien hecha que está la historia para conectar. \\ \hline
\caption{Respuestas abiertas sobre las fortalezas de \textit{Reset}.}
\label{tab:comentarios_mas_gustado}
\end{longtable}

\newpage

\section{Aspectos menos valorados (Áreas de Mejora)}
La Tabla~\ref{tab:comentarios_menos_gustado} recoge las opiniones y críticas constructivas de los usuarios al ser preguntados por el aspecto que menos les gustó del videojuego.

\begin{longtable}{|l|p{12.5cm}|}
\hline
\textbf{Usuario} & \textbf{¿Qué es lo que MENOS te ha gustado de Reset?} \\ \hline
\endfirsthead
\hline
\textbf{Usuario} & \textbf{¿Qué es lo que MENOS te ha gustado de Reset? (Continuación)} \\ \hline
\endhead
\hline
\endfoot
\hline
\endlastfoot
Usuario 1 & El nivel de \textit{Reset Bros} se complica demasiado. \\ \hline
Usuario 2 & Igual en el primer nivel hay enemigos de más para ser un primer nivel. \\ \hline
Usuario 3 & Hay algún bug pero son simples de arreglar. \\ \hline
Usuario 4 & Algún pequeño bug, pero nada que afecte a la experiencia. \\ \hline
Usuario 5 & Al guardar partida en el segundo nivel se ha reiniciado al primer nivel y no he podido avanzar. \\ \hline
Usuario 6 & Me falta feedback de los ataques que hago y recibo en el mundo 1. \\ \hline
Usuario 7 & Los enemigos amarillos en \textit{Super Reset Bros}. \\ \hline
Usuario 8 & Algunos de los enemigos en el primer mundo se vuelven frustrantes. \\ \hline
Usuario 9 & Siendo sincero, me parece que \textit{Super Reset Bros} es demasiado complicado en comparación con \textit{The Last Reset}, lo que tiene sentido, ya que debe haber una curva de dificultad para que la experiencia de juego resulte reconfortante. Y, aunque entiendo que es otro género de videojuegos muy distinto, mi sensación es que la curva de dificultad en el Nivel 2 de \textit{Super Reset Bros} es demasiado elevada para un jugador que no esté nada acostumbrado al género de las plataformas. La primera parte del Nivel 2 es sencilla y te permite familiarizarte con los controles. Pero siento que ese aprendizaje luego no sirve de tanto una vez llegas a la zona de plataformas móviles que suben, bajan e incluso que desaparecen tras pisarlas. No sé si hay muy poco tiempo en esas plataformas o si los pinchos de las paredes dan muy poco espacio para mover al personaje con libertad, pero se me ha hecho muy complicado y no he podido superarlo. Los controles están muy bien explicados en el menú, pero si el sentido es que sea difícil, está perfecto para enganchar a jugadores más experimentados. Creo que con más habilidad lo disfrutaría mucho más. \\ \hline
Usuario 10 & En general me ha encantado todo, si tuviese que decir lo que menos diría la estética de la interfaz, aún así la considero un 9/10. \\ \hline
Usuario 11 & Creo que le vendría bien un sistema de personalización de controles para facilitar la adaptación entre los dos géneros que presenta. \\ \hline
Usuario 12 & Se me ha hecho demasiado difícil el segundo nivel del mundo 2. \\ \hline
\caption{Respuestas abiertas sobre las debilidades y áreas de mejora de \textit{Reset}.}
\label{tab:comentarios_menos_gustado}
\end{longtable}
